\documentclass[11pt]{article}

\usepackage[T1]{fontenc}
\usepackage[utf8]{inputenc}
\usepackage[margin=1.02in]{geometry}
\usepackage{amsmath,amssymb,amsthm,mathtools}
\usepackage{array}
\usepackage{enumitem}
\usepackage{microtype}
\usepackage[hidelinks]{hyperref}

\newtheorem{theorem}{Theorem}
\newtheorem{lemma}{Lemma}
\newtheorem{proposition}{Proposition}
\newtheorem{definition}{Definition}
\newtheorem{claim}{Claim}
\newtheorem{remark}{Remark}
\newtheorem{corollary}{Corollary}

\newcommand{\NS}{\mathrm{NS}}
\newcommand{\Pack}{\mathrm{Pack}}
\newcommand{\Part}{\mathrm{Part}}
\newcommand{\Field}{\mathrm{Field}}
\newcommand{\Member}{\mathrm{Member}}
\newcommand{\Exit}{\mathrm{Exit}}
\newcommand{\CM}{\mathrm{CM}}
\newcommand{\Legal}{\mathrm{Legal}}
\newcommand{\GenuineCMExit}{\mathrm{GenuineCMExit}}
\newcommand{\Owork}{\mathfrak O_{\NS}^{\mathrm{work}}}
\newcommand{\WitnessRecord}{\mathcal W}
\newcommand{\Rterm}{\mathcal R_*}
\newcommand{\Ffail}{\mathcal F}
\newcommand{\Pass}{\mathrm{Pass}}
\newcommand{\Fail}{\mathrm{Fail}}
\newcolumntype{L}[1]{>{\raggedright\arraybackslash}p{#1}}
\setlength{\emergencystretch}{1em}

\title{A Class-Membership Contrapositive Route for the Navier--Stokes Millennium Problem}
\author{Thomas Birnie}
\date{Working manuscript\\May 25, 2026}

\begin{document}
\maketitle

\begin{abstract}
The usual route to the Navier--Stokes Millennium problem is forward-positive.
It tries to prove enough a priori control to continue the smooth solution forever.
That route remains the gold standard because it proves smoothness directly.
It also produces a long list of hard terminal obstructions: source concentration, zero-radius residues, retained-amplitude growth, height flux, signed current, square reserve, endpoint readout, no-pulse failure, and comparison branches.

The route developed here is class-membership contrapositive.
The same smooth-data Navier--Stokes evolution is organized as a working class object.
A legal continuation branch is a member branch.
A finite terminal branch capable of nonsmoothness is tested as a same-solution terminal witness.
The witness must retain three services: a positive same-fluid carrier, the same pressure-viscosity participation tower, and positive-scale one-field coherence.
These are denoted by \(\Pack_Q\), \(\Part_{N,Q}\), and \(\Field_{N,r,Q}\).
The terminal classification theorem shows that every same-solution terminal witness loses one of these services in first-face order.
The lost service gives class exit:
\[
  \Exit(Q;\Owork):=\neg\Member(Q;\Owork).
\]
The pass branch remains a member branch and reads out to the usual \(H^s\), \(s>5/2\), continuation theorem.
The fail branch exits the class of legal same-solution continuation candidates.
The Clay-level closing law is then isolated cleanly: smooth original data may not produce a genuine finite class exit.
This is the silver-standard route: the pass-side control mechanism may vary, while every breakdown-capable obstruction is forced to declare the exact class service it has lost.
\end{abstract}

\setcounter{tocdepth}{1}
\tableofcontents
\clearpage

\section{The Problem And The Change Of Question}

The incompressible three-dimensional Navier--Stokes equations are
\[
  \partial_t u +(u\cdot\nabla)u+\nabla p=\nu\Delta u,\qquad
  \nabla\cdot u=0,
\]
with smooth divergence-free initial data \(u_0\).
The Clay problem asks whether the corresponding smooth solution exists for all time.

The direct proof question is:
\[
  \text{which estimate forces the smooth solution to continue?}
\]
That question is natural, and it is the forward-positive gold standard.
One proves a bound strong enough for classical continuation, usually through a continuation norm such as \(H^s\), \(s>5/2\).

The route here changes the question after the forward proof reaches a genuine obstruction.
It asks:
\[
  \begin{gathered}
  \text{can a branch capable of nonsmooth terminal failure}\\
  \text{remain a legal branch of the same class object?}
  \end{gathered}
\]
The two questions have different burdens.
The forward question tries to erase the bad object.
The class-membership question lets the bad object appear and asks whether it can still be the same kind of object the Clay problem asks us to continue.

The basic split is this.
Let \(O\) be a same-solution obstruction reached by a forward Navier--Stokes proof.
The obstruction has a pass side and a fail side:
\[
  O_{\Pass}\Rightarrow \text{the same continuation services remain},
\]
\[
  O_{\Fail}\Rightarrow \text{the branch is capable of finite terminal nonsmoothness}.
\]
The pass side is the member branch.
The fail side is tested as a terminal witness and must lose one of the required services.
Once it loses a required service, it exits the class:
\[
  O_{\Fail}\Rightarrow \Exit(Q;\Owork).
\]
Thus the class proof separates two mathematical jobs.
First, the fail side must enter the terminal witness tree and leave the class by a first failed service.
Second, the original smooth-data evolution must have no genuine finite class exit.
The first job is the terminal classification theorem.
The second job is the Clay closure law.

\section{The Forward Obstructions}

The forward-positive program remains the right way to find the hard objects.
It asks each obstruction to be controlled directly.
The present proof uses those obstructions as terminal test objects.

\subsection{Source Control}

Source control tries to prevent terminal concentration of the nonlinear source tree.
Its strongest form is a scale-critical Carleson-type control on donor refill, source reserve, and active ancestry.
The obstruction is that first-moment or local-energy control can leave a square or scale-critical reserve unpaid.

In the class-membership route, selected unpaid donor ancestry is a carrier question first.
A branch whose terminal ancestry has no positive same-fluid carrier fails \(\Pack_Q\).
A retained source residue with a carrier then faces the participation question.
Only after carrier and participation survive may a retained source residue become a Field-readout question.

\subsection{Zero-Radius And Zeno Objects}

Zeno language describes terminal atoms, zero heat-time residues, temporal non-atomicity, and rigid residue classes.
A forward proof tries to rule out those objects positively.
The class-membership proof asks where the selected terminal object lives.

A zero-radius terminal object has no positive terminal carrier.
It fails \(\Pack_Q\).
A retained positive-scale Zeno object stays inside the same Pack/Part/Field witness tree and is read by first failed service.
Zeno is a presentation of terminal behaviour inside the three-service tree.

\subsection{Retained Amplitude And Low Strain}

Retained-amplitude arguments try to turn surviving carrier and participation structure into enough amplitude control for endpoint coherence.
The obstruction is that bounded deformation data can coexist with accumulated positive low-strain action.
Height-flux and low-strain routes sharpen the same pressure point: near-band source charge and frame variation return to the signed source balance.

The class-membership route reads this family as a finite routing engine.
The selected branch lands as Pack gain failure, Part participation or dwell failure, legal retained Field failure, or a Zeno subcase already sorted by the witness tree.

\subsection{Signed Current}

Signed-current arguments try to rectify terminal source-current motion into a no-free-sink or parabolic edge-resistance law.
The obstruction is the variable sign of the strain-production term and the scale-critical flux-amplitude remainder.
Those terms prevent a free coercion theorem.

For class membership, signed current has no independent face.
It returns to Pack, Part, retained Field after licensing, Zeno, or positive supplier support.
Its value is that it identifies where the same services fail.

\subsection{Square Reserve}

Square-reserve arguments ask for a square-source charge law strong enough to control terminal reserve.
They matter because a first-moment source tail can vanish while a square reserve remains active.
The forward branch seeks a direct evolution or charge theorem.

The class-membership branch asks what the selected reserve failure is.
Detached reserve ancestry is Pack-side.
A retained reserve with carrier and participation can become a Field diagnostic.
An absent square-charge identity, by itself, is supplier work.

\subsection{Endpoint Readout And No-Pulse}

Endpoint readout asks for a good scale, a read cover, a coherent field, and an \(H^s\) continuation bound.
No-pulse and averaged routes try to prevent terminal pulses, averaged jumps, and terminal-tail escape.
These are close to classical continuation, so they can sound like the whole proof.

They are downstream in the class-membership argument.
Readout explains why a surviving Pack/Part/Field packet continues.
No-pulse can support retained Field or endpoint production.
Neither replaces the primitive class test.

\subsection{Local Critical Translators And Comparison Branches}

Local \(L^3\) and Duhamel translators are useful when a public-critical-looking terminal mass must be brought back to the same witness record.
They turn a local packet into native source work or a Pack/Part/Field exit.
Their job is branch-local.

Fixed-viscosity Euler-to-Navier--Stokes transfer and Euler mirror comparisons serve different theorem programs.
They illuminate where viscosity enters and where inviscid comparison diverges.
They enter this proof only through a theorem that lands the exact Navier--Stokes object in the same Pack/Part/Field class test.

\section{The Working Class Object}

The proof needs a class object because the contrapositive is a class-law argument.
The class is built from the same datum, same equation, same pressure, same viscosity, and same fluid branch.

\begin{definition}[Working Clay class object]
The working Clay class object \(\Owork\) is the class of same-solution Navier--Stokes continuation packets generated by the original smooth datum and read on one terminal witness record.
For a terminal packet \(Q\), the predicate
\[
  \Member(Q;\Owork)
\]
means that \(Q\) retains the services needed to be a legal same-solution continuation candidate.
The class-exit predicate is
\[
  \Exit(Q;\Owork):=\neg\Member(Q;\Owork).
\]
\end{definition}

\begin{definition}[Witness record]
A witness record is the common record
\[
  \WitnessRecord=(u_0,\nu,u,p,Q,N,r,\Phi,T_*),
\]
where \(u_0\) is the original datum, \((u,p)\) is the same Navier--Stokes solution on \([0,T_*)\), \(Q\) is the terminal packet, \(N\) is the participation depth, \(r\) is the positive field scale, \(\Phi\) is the same-fluid carrier family, and \(T_*\) is the alleged terminal time.
\end{definition}

The common-record condition is important.
A carrier, participation tower, and readout synthesize into one continuation packet only when they belong to the same terminal object.
The class test is a single-record test.

\begin{definition}[The three services]
On a witness record \(\WitnessRecord\), the services are:
\[
  \Pack_Q,\qquad \Part_{N,Q},\qquad \Field_{N,r,Q}.
\]
\(\Pack_Q\) is the positive same-fluid carrier service.
It means that there is a finite transported carrier family \(\{\Phi_a\}_{a=1}^m\), a finite-overlap cover \(\{C_a\}_{a=1}^m\), and a number \(\rho_Q>0\) such that the selected terminal packet \(Q\) is carried on sets of diameter at least \(\rho_Q\), all transported by the same velocity field \(u\).

\(\Part_{N,Q}\) is the same-equation participation service.
It means that on the carrier supplied by \(\Pack_Q\), the differentiated pressure-viscosity tower remains the same Navier--Stokes tower through depth \(N\):
\[
  D_tU_k=K_k+B_k,\qquad
  K_k=-\nabla^{k+1}p+\nu\Delta U_k,\qquad 0\le k\le N.
\]
The tower has one finite participation envelope on the witness record:
\[
  \mathcal T_{N,Q}
  :=
  \sum_{k=0}^{N}
  \sup_{t<T_*}
  \sum_{a=1}^{m}
  \|\chi_a(t)U_k(t)\|_{L^2}^2
  <\infty,
\]
where \(\{\chi_a\}\) is a subordinate partition for the carrier cover.

\(\Field_{N,r,Q}\) is the positive-scale one-field coherence service.
It means that after carrier and participation survive, one positive scale \(0<r\le \rho_Q\) carries a finite coherent field energy:
\[
  \mathcal E_{N,r,Q}
  :=
  \sup_{t<T_*}
  \sum_{a=1}^{m}
  \sum_{|\alpha|\le N}
  \|\chi_a(t)\nabla^\alpha u(t)\|_{L^2}^2
  <\infty,
\]
with compatible pressure readout on the same cover.
\end{definition}

\begin{definition}[CM packet]
The class-membership packet at depth \(N\), scale \(r\), and terminal packet \(Q\) is
\[
  \CM_{N,r,Q}
  :=
  \Pack_Q\wedge\Part_{N,Q}\wedge\Field_{N,r,Q}
\]
on one witness record.
\end{definition}

\begin{lemma}[Membership readout]
Let \(N_s\) be a depth high enough for the standard \(H^s\), \(s>5/2\), continuation criterion.
If a legal same-solution terminal packet remains a member of \(\Owork\), then
\[
  \Member(Q;\Owork)
  \Longrightarrow
  \sup_{t<T_*}\|u(t)\|_{H^s}<\infty
\]
for some \(s>5/2\), and the same Navier--Stokes solution extends classically beyond \(T_*\).
\end{lemma}

\begin{proof}
Membership means that the terminal packet retains the three services on one witness record.
\(\Pack_Q\) supplies a finite positive carrier cover with finite overlap.
\(\Part_{N_s,Q}\) identifies the differentiated fields on that cover with the same pressure-viscosity Navier--Stokes tower through the continuation depth.
\(\Field_{N_s,r,Q}\) supplies the finite local derivative energy on a positive scale.
The partition of unity subordinate to the finite cover sums those local bounds into the displayed \(H^s\) bound at a supercritical continuation level \(s>5/2\).
Classical local well-posedness then continues the same solution past \(T_*\).
\end{proof}

\section{Which Objects Are Tested}

The class-membership proof tests finite terminal objects offered as Clay breakdown witnesses for the same original Navier--Stokes solution.

\begin{definition}[CM-test-admissible Clay terminal object]
A finite terminal object \(W_*\) is CM-test-admissible when it is offered as a breakdown witness for the same original smooth-data Navier--Stokes solution.
This means:
\begin{enumerate}[leftmargin=2em]
\item \(W_*\) is attached to the maximal classical solution generated by the original datum \(u_0\).
\item On every preterminal window, the object solves the original incompressible Navier--Stokes system with the same velocity and pressure.
\item The terminal extraction uses the same transported fluid family supplied by the preterminal smooth windows.
\item After legal boundary losses, cutoff losses, non-selected branches, readout-only residues, and paid finite terms are removed, \(W_*\) is still the object claiming to be the finite terminal obstruction.
\item The remaining terminal services available to the object are exactly \(\Pack_Q\), \(\Part_{N,Q}\), and \(\Field_{N,r,Q}\).
\end{enumerate}
\end{definition}

\begin{remark}
CM-test-admissible means admissible for the test.
Membership and exit are conclusions.
\end{remark}

\begin{lemma}[Finite Clay breakdown enters the test]
Every finite same-solution Clay breakdown witness is CM-test-admissible.
\end{lemma}

\begin{proof}
A finite same-solution Clay breakdown witness is, by definition, the original smooth-data Navier--Stokes solution stopped at a finite terminal time where classical continuation is alleged to fail.
The datum, equation, pressure, velocity, viscosity, and preterminal fluid family are therefore the same ones used before the terminal time.
Legal losses and irrelevant branches can be stripped away without changing the alleged obstruction.
The remaining object is precisely the terminal object to be tested by the three services.
\end{proof}

\section{Terminal Packet Capture}

The capture theorem is the bridge from a hypothetical finite breakdown to the Pack-first decision tree.

\begin{theorem}[Terminal packet capture]
Let \((u,p)\) be the maximal smooth Navier--Stokes solution generated by \(u_0\) on \([0,T_*)\), with \(T_*<\infty\) assumed.
Let \(\Rterm\) be a same-solution terminal obstruction after legal losses, non-selected branches, readout-only residues, and paid finite terms have been removed.
Fix \(N_s\) high enough for \(H^s\), \(s>5/2\), continuation readout.
Then
\[
  \Rterm
  \Longrightarrow
  \neg\Pack_Q
  \vee
  \neg\Part_{N_s,Q}
  \vee
  \forall r>0\,\neg\Field_{N_s,r,Q}.
\]
Equivalently,
\[
  \Rterm
  \Longrightarrow
  \begin{cases}
  \neg\Pack_Q,\\
  \Pack_Q\wedge\neg\Part_{N_s,Q},\\
  \Pack_Q\wedge\Part_{N_s,Q}\wedge\forall r>0\,\neg\Field_{N_s,r,Q}.
  \end{cases}
\]
\end{theorem}

\begin{proof}
For every \(\varepsilon>0\), the restriction of \((u,p)\) to \([0,T_*-\varepsilon]\) is classical.
Each compact preterminal window carries the same datum, same equation, same pressure, same viscosity, same transported fluid family, and hence the same local continuation services.

Pass to a terminal subsequence producing \(\Rterm\).
The selected object first asks whether a positive same-fluid carrier survives.
If no such carrier survives, \(\neg\Pack_Q\).

Assume that \(\Pack_Q\) survives.
The next question is whether the same differentiated Navier--Stokes pressure-viscosity tower survives on that carrier through depth \(N_s\).
If the tower fails, \(\neg\Part_{N_s,Q}\).

Assume that both \(\Pack_Q\) and \(\Part_{N_s,Q}\) survive.
If there is a positive scale \(r>0\) with \(\Field_{N_s,r,Q}\), then the full packet \(\CM_{N_s,r,Q}\) survives on the same witness record.
By the membership readout lemma, the same solution has a bounded \(H^s\), \(s>5/2\), continuation norm and extends classically past \(T_*\).
That contradicts the terminal role of \(\Rterm\).
Thus \(\Field_{N_s,r,Q}\) fails for every \(r>0\).
\end{proof}

\begin{corollary}[No fourth primitive service]
After legal losses, readout-only residues, non-selected branches, and paid finite terms are removed, a same-solution terminal obstruction has no primitive service outside Pack, Part, and Field.
\end{corollary}

\begin{proof}
The preceding theorem tests exactly the services needed for the same-solution continuation packet.
Any proposed residue either lacks the carrier, lacks participation after carrier survives, lacks positive-scale coherence after carrier and participation survive, or is a supplier/readout description waiting for one of those services to be tested.
\end{proof}

\section{Branch-Family Exhaustion}

The previous theorem explains how the obstruction families enter the proof.
Each family enters through the same question: what terminal object is selected, and which service fails first?

\begin{proposition}[Branch-family exhaustion]
Every non-Euler Navier--Stokes branch family used by the working class-membership program has one of the following mathematical roles:
\[
  \begin{gathered}
  \Pack\text{-exit},\qquad
  \Part\text{-exit},\qquad
  \Field\text{-exit},\\
  \text{finite typed disjunction},\qquad
  \text{auxiliary until a named face theorem lands it}.
  \end{gathered}
\]
\end{proposition}

\begin{proof}
We read the families by their selected terminal object.

Source-control and donor-refill branches select terminal source ancestry or reserve.
Detached or unpaid ancestry without a positive same-fluid carrier is Pack-exit.
Retained source residue with carrier and broken pressure-viscosity participation is Part-exit.
Retained source residue that has passed Pack and Part becomes a Field/readout diagnostic.

Zero-radius and terminal-atom branches select an object with no positive terminal carrier unless a retained positive-scale packet has first been supplied.
The zero-radius branch is Pack-exit.
Retained Zeno alternatives are finite typed subcases within the same witness tree.

Retained-amplitude, height-flux, low-strain, signed-current, and square-reserve branches select failures of gain, dwell, participation, source balance, or retained coherence.
Each such failure either breaks the carrier, breaks the same pressure-viscosity participation tower after the carrier survives, breaks positive-scale coherence after Pack and Part survive, or remains a positive supplier theorem outside the class primitive.

Good-scale, endpoint-readout, no-pulse, averaged, and terminal-tail branches read a packet that has already retained enough structure to be read.
They support the Field and continuation consequences.
They become proof material only when the exact packet has already entered Pack, Part, or Field.

Local critical translators place local \(L^3\) or Duhamel mass back on the same witness record.
Their role is branch-local: native source work or Pack/Part/Field exit.

Fixed-viscosity transfer and Euler mirror branches are comparison programs.
They remain outside the non-Euler class-membership proof until a theorem transfers the exact Navier--Stokes object into the same witness tree.

These alternatives exhaust the branch families because every branch either supplies a selected same-solution terminal object or supplies context for such an object.
Selected terminal objects are exhausted by terminal packet capture.
Context without a selected face remains support.
\end{proof}

\begin{center}
\small
\begin{tabular}{L{0.30\linewidth}L{0.31\linewidth}L{0.30\linewidth}}
\textbf{Branch family} & \textbf{Forward-positive burden} & \textbf{CM role}\\
\hline
Source control and donor refill & prove scale-critical source reserve control & Pack first; Part after carrier survives; retained Field only after Pack and Part\\
Zero-radius and Zeno & prove temporal non-atomicity or rigid residue exclusion & zero-radius Pack exit; retained positive-scale typed subcases\\
Retained amplitude and low strain & prove gain, BV, or height-flux control & Pack/Part/Field/Zeno finite routing\\
Signed current & prove no-free-sink or edge resistance & auxiliary route returning to Pack, Part, retained Field, or Zeno\\
Square reserve & prove square-source charge or native trilinear control & auxiliary theorem unless selected reserve lands in first-face order\\
Endpoint readout and no-pulse & prove good scale, read cover, and pulse exclusion & downstream support for membership readout and Field\\
Local critical translators & localize public-critical terminal mass & branch-local native source work or class exit\\
Transfer and mirror comparison & compare neighbouring theorem worlds & outside this proof until an exact transfer lands\\
\end{tabular}
\end{center}

\section{Representation Standard}

The proof material is large because the forward-positive search produces many versions of the same few terminal pressures.
The manuscript is short only when the compression is mathematical.
Every source object must enter by proof role.
It must become a formal class law, a typed terminal branch, support for a named hinge, a comparison pressure-test, or a live obstruction to one of the closing hinges.

\begin{center}
\begingroup
\setlength{\tabcolsep}{3pt}
\scriptsize
\begin{tabular}{L{0.18\linewidth}L{0.18\linewidth}L{0.32\linewidth}L{0.20\linewidth}}
\textbf{Material} & \textbf{Entry in the manuscript} & \textbf{Required standard} & \textbf{Failure signal}\\
\hline
Core definitions and class laws & formal definitions, lemmas, and theorems & The datum, equation, pressure, viscosity, witness record, membership predicate, and exit predicate are stated inside the argument. & The term is vocabulary rather than a proved object.\\
Terminal classifications & first-face theorem or proof paragraph & One selected terminal packet is tested in order: \(\Pack_Q\), then \(\Part_{N,Q}\), then \(\Field_{N,r,Q}\). & The branch has no declared first failed service.\\
Forward obstruction families & branch-family compression & The obstruction is split into pass and fail sides, with the fail side offered as a terminal witness and the pass side read as membership support. & The family remains thematic pressure without a selected terminal object.\\
Auxiliary estimates and endpoint readouts & hinge support & The estimate names the exact hinge it supports: packet capture, membership readout, or no genuine finite class exit. & The estimate strengthens intuition without moving a hinge.\\
Transfer and comparison material & comparison pressure-test & An exact Navier--Stokes object is transferred into the same Pack/Part/Field witness tree before it enters the proof spine. & The comparison remains outside the class-law argument.\\
Failed attempts and adversarial audits & live burden or pressure test & A failed attempt becomes main-text material only by attacking one of the three referee hinges. & The attempt stays as background pressure.\\
Active blockers & named theorem obligation & The blocker is stated as a mathematical obligation with its service face and closing hinge identified. & The manuscript hides an open burden inside prose.\\
\end{tabular}
\endgroup
\end{center}

This standard gives the page count its mathematical meaning.
Coverage is by proof role.
The many working records are represented when their mathematical jobs have been absorbed into the class law, the branch-family table, the Pack/Part/Field witness tree, the comparison boundary, or the referee hinges.
Material with no such job stays outside the main argument.

\section{The Obstruction Census}

The forward-positive search is useful because it has already exposed the obstruction field.
The class proof uses that search as a finite census.
Each record is read as a candidate terminal pressure on the same solution, then placed into one of five roles:
\[
  \begin{gathered}
  \Pack\text{-face},\qquad
  \Part\text{-face},\qquad
  \Field\text{-face},\\
  \text{finite face disjunction},\qquad
  \text{support awaiting a selected face}.
  \end{gathered}
\]
This is the reason the manuscript can be expanded without becoming a catalogue.
The expansion must show the mathematical decision rule for each family.
The right unit is one paragraph per proof role.

\begin{theorem}[Non-Euler obstruction census]
The non-Euler Navier--Stokes working field contains \(1944\) classified proof records.
After authority carriers are separated, \(1359\) records are certified as class-exit records through Pack, Part, or Field; \(44\) inherit an already installed face rule and introduce no new selected failure; \(541\) are support records barred from the class proof until a named Pack/Part/Field landing is supplied.
The face counts are:
\[
  \Pack_Q:\ 340,\qquad
  \Part_{N,Q}:\ 346,\qquad
  \Field_{N,r,Q}:\ 865,
\]
with \(585\) records carrying no independent class face.
\end{theorem}

\begin{proof}
The census is performed on non-Euler Navier--Stokes material only.
Each record is first tested for a selected same-solution terminal failure.
When the record itself is an authority carrier, it contributes no new failure and inherits the governing face rule.
When the record selects a same-solution terminal failure, the first available failure face is read in Pack-first order.
When the record contains a supplier estimate, readout estimate, comparison branch, historical note, or build/provenance material without a selected same-solution failure, it is barred from the class proof until a theorem lands its exact object in Pack, Part, or Field.
Counting those alternatives gives the stated totals.
\end{proof}

\begin{corollary}[What the census proves]
For every non-Euler record with a selected same-solution terminal failure, the record supports one of the following conclusions:
\[
  \neg\Pack_Q,\qquad
  \Pack_Q\wedge\neg\Part_{N,Q},\qquad
  \Pack_Q\wedge\Part_{N,Q}\wedge\forall r>0\,\neg\Field_{N,r,Q},
\]
or else a finite disjunction of these three outcomes.
\end{corollary}

\begin{proof}
This is terminal packet capture read record by record.
A selected terminal object either lacks the positive carrier, loses the same-equation participation service after the carrier survives, loses positive-scale one-field coherence after carrier and participation survive, or has not yet been promoted into the class test.
The promoted cases are exactly the Pack, Part, Field, and finite-disjunction cases.
\end{proof}

\section{The Pack Face}

The Pack face asks the first physical question.
Does the alleged terminal object retain a positive same-fluid carrier from the original smooth solution?
A branch that cannot answer yes has already left the class.
It does not matter how much formal source language remains around it.
Without a positive carrier, there is no same-solution packet to continue.

\begin{definition}[Pack-side terminal failure]
A selected terminal object has Pack-side failure when no positive transported carrier family survives for the packet \(Q\):
\[
  \neg\Pack_Q.
\]
The typical Pack-side mechanisms are zero terminal radius, detached donor ancestry, terminal source atom without positive carrier, unsupported first-created square reserve, and terminal carrier birth with no charge path.
\end{definition}

\begin{lemma}[Zero-radius packets are Pack-side]
Let \(Q\) be a selected terminal packet whose transported support has no positive terminal radius:
\[
  \rho_Q=0.
\]
Then \(Q\) fails \(\Pack_Q\) and hence exits \(\Owork\).
\end{lemma}

\begin{proof}
\(\Pack_Q\) requires a finite transported carrier cover with a positive lower carrier scale.
The condition \(\rho_Q=0\) removes that service.
The object may still be visible as a terminal trace or residue, but it is not a same-fluid continuation packet.
Thus \(\neg\Pack_Q\), and class exit follows by the definition of \(\Exit\).
\end{proof}

\begin{lemma}[Detached donor ancestry is Pack-side]
Suppose a terminal source or reserve branch selects donor ancestry that is not represented on a positive same-fluid carrier for \(Q\).
Then the selected branch is Pack-side:
\[
  \text{detached ancestry}(Q)\Longrightarrow\neg\Pack_Q.
\]
\end{lemma}

\begin{proof}
Donor ancestry can support a terminal obstruction only when the ancestry is still carried by the same transported fluid family.
Detached ancestry has no such carrier.
It may indicate where a forward-positive source estimate failed, but it does not supply a legal same-solution terminal packet.
The first service fails at Pack.
\end{proof}

\begin{lemma}[First-created reserve must be paid or carried]
Let \(R_N(W)\) be a positive square reserve on an admissible terminal heat window \(W\).
If the positive first-created increment is neither inherited from the parent window nor paid by a charge on \(W\), then the selected reserve branch is not a member branch.
It is either a Pack-side carrier birth failure or a live Pack-support burden.
\end{lemma}

\begin{proof}
A member branch cannot acquire positive terminal reserve from nowhere.
Positive reserve appearing at the terminal window must be represented by inherited mass, declared loss, or a charge path tied to the same transported packet.
When no such representation is available, the candidate branch is not yet a member packet.
At the class level the branch is Pack-side unless a later theorem supplies the missing positive carrier and charge path.
\end{proof}

The current Pack-side live burden is the source-reserve birth chain.
The useful form is not a global smoothness estimate.
The useful form is a same-packet creation law:
\[
  R_N(W)
  \le
  (1-\delta)R_N(\mathrm{Past}(W))
  +\mathrm{Loss}_N(W)
  +\mathrm{Charge}_N(W)
  +o_N(1).
\]
This law says that terminal reserve cannot be first-created for free.
The latest reduction sharpens the remaining positive source remainder into a source-weighted angular depletion burden on the selected positive source carrier.
In the manuscript this belongs at the no-genuine-exit hinge, not at membership readout.

\section{The Part Face}

The Part face asks whether the selected carrier still participates in the same Navier--Stokes pressure-viscosity tower.
Pack alone is only a carrier.
The Clay solution is not a moving set; it is the same velocity-pressure solution of the same equation.
The Part face protects that identity.

\begin{definition}[Part-side terminal failure]
Assume \(\Pack_Q\) survives.
A selected terminal object has Part-side failure when the differentiated Navier--Stokes tower cannot be retained on the carrier through the continuation depth:
\[
  \Pack_Q\wedge\neg\Part_{N,Q}.
\]
The typical Part mechanisms are pressure participation loss, viscosity participation loss, dwell failure, same-object dynamic carrier failure, and incompatible local source participation.
\end{definition}

\begin{lemma}[Pressure-viscosity loss is Part-side]
Let \(Q\) retain a positive carrier.
If, on that carrier, the differentiated fields cannot be represented by the same pressure-viscosity tower
\[
  D_tU_k=-\nabla^{k+1}p+\nu\Delta U_k+B_k,\qquad 0\le k\le N,
\]
then \(Q\) fails \(\Part_{N,Q}\).
\end{lemma}

\begin{proof}
The Part service is exactly the service that keeps the carrier tied to the original Navier--Stokes equation.
When the pressure, viscosity, or nonlinear participation terms cease to be the same tower on the selected carrier, the object may still be a transported object, but it is no longer a legal same-solution continuation packet.
Therefore the first surviving failure after Pack is \(\neg\Part_{N,Q}\).
\end{proof}

\begin{lemma}[Local translators are Part tests before they are readouts]
A local critical packet, local \(L^3\) packet, or Duhamel packet enters the class proof only after it is placed on the same witness record.
Before that placement it is support.
After that placement it is tested first for Pack and then for Part.
\end{lemma}

\begin{proof}
Local critical information can describe a terminal concentration without identifying the same witness record.
The class proof needs the same datum, same pressure, same viscosity, same carrier family, and same terminal packet.
The translator's job is to recover that common record.
Once the record is recovered, the packet faces the ordinary Pack and Part tests.
Only after those services survive can the packet become a Field or endpoint readout object.
\end{proof}

\begin{proposition}[Part is the same-equation guard]
A branch with a carrier and a broken pressure-viscosity tower cannot be a legal continuation branch for the Clay problem.
\end{proposition}

\begin{proof}
A legal continuation branch must solve the same incompressible Navier--Stokes equations with the same pressure and viscosity.
The carrier alone does not supply that identity.
When the pressure-viscosity tower breaks, the selected object is no longer the same equation on the same packet.
Thus it exits before any endpoint readout is relevant.
\end{proof}

\section{The Field Face}

The Field face is reached only after Pack and Part survive.
At that point the selected object has a positive same-fluid carrier and the same pressure-viscosity participation tower.
The remaining question is whether there is one positive scale carrying coherent field data strong enough for continuation.

\begin{definition}[Field-side terminal failure]
Assume \(\Pack_Q\) and \(\Part_{N,Q}\) survive.
The packet has Field-side failure when every positive scale loses coherent field energy:
\[
  \Pack_Q\wedge\Part_{N,Q}\wedge\forall r>0\,\neg\Field_{N,r,Q}.
\]
The typical Field mechanisms are retained-amplitude blowup, terminal pulse, averaged jump, good-scale failure, and incompatible surviving traces.
\end{definition}

\begin{lemma}[Retained positive-window blowup is Field-side]
Suppose \(\Pack_Q\) and \(\Part_{N,Q}\) survive.
If a retained positive terminal window has no positive scale on which the \(N\)-level field energy remains finite, then the selected branch fails Field.
\end{lemma}

\begin{proof}
After Pack and Part survive, the only service left for membership readout is the positive-scale coherent field service.
Unbounded retained field energy on every positive scale removes that service.
The branch is therefore Field-side.
\end{proof}

\begin{lemma}[Endpoint readout is downstream of Field]
Good-scale, no-pulse, terminal-tail, averaged, and endpoint-cover arguments are continuation readout support.
They prove membership only after the packet has already retained Pack, Part, and Field on one witness record.
\end{lemma}

\begin{proof}
Endpoint readout uses the packet's existing services.
It does not create the positive carrier, the same pressure-viscosity tower, or the coherent field service by itself.
When those services survive, readout converts them into an \(H^s\), \(s>5/2\), continuation bound.
When one service fails, readout is no substitute for the missing service.
\end{proof}

\begin{proposition}[Field is the last continuation gate]
A same-solution terminal packet that retains Pack, Part, and Field at a continuation depth \(N_s\) cannot be a finite nonsmooth terminal obstruction.
\end{proposition}

\begin{proof}
Pack gives the finite positive carrier cover.
Part keeps the same pressure-viscosity tower on that cover.
Field gives finite coherent derivative energy on a positive scale.
The finite cover and partition of unity produce a bounded \(H^s\), \(s>5/2\), norm.
The classical continuation theorem extends the same solution past the alleged terminal time.
Thus a terminal obstruction must fail before or at Field.
\end{proof}

\section{The Support Boundary}

Many strong estimates do real work without being class faces.
This is where the manuscript must be careful.
A source-control estimate, no-pulse estimate, signed-current estimate, or comparison estimate may be essential to a route.
It still has to say which class service it protects.

\begin{proposition}[Support cannot replace a face]
Let \(S\) be a source, readout, comparison, or supplier statement.
If \(S\) does not land a selected same-solution object in \(\Pack_Q\), \(\Part_{N,Q}\), \(\Field_{N,r,Q}\), membership readout, terminal packet capture, or no genuine finite exit, then \(S\) is support for the class proof and not a class-proof conclusion.
\end{proposition}

\begin{proof}
The class proof has only three witness services and three closing hinges.
A supplier estimate can help prove one of them, and a readout estimate can explain the continuation consequence after the services survive.
A comparison estimate can show why a neighbouring theorem world behaves differently.
None of those roles is itself membership or exit.
The statement enters the proof only when its exact target is a class service or closing hinge.
\end{proof}

\begin{corollary}[Forward-positive estimates keep their value]
Demoting a forward-positive estimate from class authority does not discard it.
It assigns the estimate a proof job:
\[
  \text{support Pack},\quad
  \text{support Part},\quad
  \text{support Field},\quad
  \text{support readout},\quad
  \text{or support no genuine exit}.
\]
\end{corollary}

\begin{proof}
The estimate remains available exactly where it has mathematical force.
The class proof only prevents the estimate from pretending to be the class law by itself.
Once the estimate names a service or hinge, it can be used in the proof of that service or hinge.
\end{proof}

\section{Comparison And Transfer Boundary}

The Navier--Stokes class proof is not an Euler proof and not a fixed-viscosity transfer proof.
Euler comparison can be illuminating because it shows what disappears when viscosity and parabolic participation are absent.
It becomes Navier--Stokes proof material only by landing an exact Navier--Stokes packet in the same witness tree.

\begin{proposition}[Comparison boundary]
Let \(C\) be an Euler, fixed-viscosity transfer, or neighbouring-equation comparison branch.
The branch enters the present Navier--Stokes class proof only through a theorem of the form
\[
  C\Longrightarrow
  \text{selected Navier--Stokes packet }Q
  \text{ has a Pack, Part, or Field decision}.
\]
Before that theorem, \(C\) is comparison pressure.
\end{proposition}

\begin{proof}
The working class object \(\Owork\) is built from the original Navier--Stokes datum, pressure, viscosity, and same-fluid participation.
Comparison branches may share shapes with Navier--Stokes branches, but the class proof needs the exact same-solution packet.
The transfer theorem supplies that exact packet.
Without it, the comparison branch has no class face inside \(\Owork\).
\end{proof}

\section{The Remaining Closing Burden}

The class-law proof separates terminal classification from final Clay closure.
The classification theorem says that every finite same-solution failure witness exits by Pack, Part, or Field.
The final Clay step says that the original smooth-data evolution never produces a genuine finite class exit.
That is the live mathematical burden:
\[
  \text{OriginalSmoothData}(u_0)
  \Longrightarrow
  \neg\GenuineCMExit(T_*;u_0)
  \qquad(T_*<\infty).
\]

\begin{definition}[No-genuine-exit burden]
The no-genuine-exit burden is the theorem that original smooth Navier--Stokes data cannot generate a same-solution terminal packet that is CM-test-admissible and outside \(\Owork\).
It is a class-law theorem about the original branch, not a request to erase every possible bad object in an unrelated class.
\end{definition}

\begin{proposition}[Current Pack-side reduction]
The present no-genuine-exit pressure is Pack-side.
It asks for a source-reserve creation and depletion law strong enough to prevent a retained positive-scale native source carrier from appearing as unpaid terminal reserve.
\end{proposition}

\begin{proof}
The active source-reserve branch concerns first-created positive source reserve on terminal heat windows.
Such a branch is a carrier and source-birth question before it is a field readout question.
The current reductions place the work on reserve creation, terminal first appearance, positive remainder depletion, and source-weighted angular depletion.
Those are Pack-support obligations for the no-genuine-exit hinge.
\end{proof}

\begin{theorem}[Completion criterion]
If terminal packet capture, membership readout, branch-family exhaustion, and the no-genuine-exit burden all hold, then the smooth Navier--Stokes solution generated by smooth divergence-free data exists for all time.
\end{theorem}

\begin{proof}
Terminal packet capture and branch-family exhaustion classify every finite same-solution terminal failure as Pack, Part, Field, finite disjunction, or unlicensed support.
Unlicensed support has no selected class failure and cannot be a terminal counterexample.
Membership readout continues every in-class branch.
The no-genuine-exit burden removes every out-of-class branch produced by the original smooth data.
Thus a finite terminal time has no legal branch left.
The solution continues globally.
\end{proof}

\section{Class Exit}

The class-exit theorem turns first-face failure into the contrapositive conclusion.

\begin{theorem}[Finite failure witness exits the class]
Let \(W_*\) be a CM-test-admissible finite Clay terminal object.
Then for a suitable terminal packet \(Q\) and continuation depth \(N_s\),
\[
  W_*
  \Longrightarrow
  \Exit(Q;\Owork).
\]
\end{theorem}

\begin{proof}
By terminal packet capture, \(W_*\) loses a first required service:
\[
  \neg\Pack_Q
  \vee
  \neg\Part_{N_s,Q}
  \vee
  \forall r>0\,\neg\Field_{N_s,r,Q}.
\]
If \(\Pack_Q\) fails, the terminal object has no positive same-fluid carrier and is outside the member packets.
If \(\Part_{N_s,Q}\) fails after Pack survives, the terminal object has left the same pressure-viscosity Navier--Stokes tower and is outside the member packets.
If every positive \(\Field_{N_s,r,Q}\) fails after Pack and Part survive, the terminal object lacks the positive-scale coherence required for membership readout and is outside the member packets.

Each case is the loss of a required service for \(\Member(Q;\Owork)\).
Therefore
\[
  \Exit(Q;\Owork):=\neg\Member(Q;\Owork).
\]
\end{proof}

\section{The Class-Law Proof}

Classification supplies the exit witness.
The remaining question is whether a finite exit branch can still be a legal continuation of the Clay solution from smooth data.

\begin{definition}[Genuine finite CM exit]
For the maximal smooth solution generated by \(u_0\) on \([0,T_*)\), with \(T_*<\infty\), write
\[
  \GenuineCMExit(T_*;u_0)
\]
when a same-solution terminal packet \(Q\) is CM-test-admissible and satisfies \(\Exit(Q;\Owork)\).
This is a terminal class exit produced by the original smooth-data evolution itself.
\end{definition}

\begin{proposition}[Legal branch law]
Every legal same-solution continuation branch for the Clay Navier--Stokes problem is represented by the working class object:
\[
  \Legal_{\NS}(Q)\Longrightarrow \Member(Q;\Owork).
\]
\end{proposition}

\begin{proof}
A legal same-solution continuation branch is the original smooth-data Navier--Stokes evolution continued as the same velocity-pressure solution with the same viscosity and same fluid participation.
To be such a continuation branch, it must retain a positive same-fluid carrier, the same pressure-viscosity participation tower through the continuation depth, and positive-scale coherent field data for the \(H^s\), \(s>5/2\), readout.
Those are exactly the services defining membership in \(\Owork\).
\end{proof}

\begin{theorem}[No third in-class terminal branch]
There is no terminal branch that is both an in-class member branch and a finite nonsmooth obstruction branch:
\[
  \Member(Q;\Owork)
  \quad\text{and}\quad
  Q\text{ blocks smooth continuation at }T_*<\infty.
\]
\end{theorem}

\begin{proof}
Membership gives the three services on one witness record.
By membership readout, those services give a bounded \(H^s\), \(s>5/2\), continuation norm.
Classical local theory then extends the same solution past \(T_*\).
That contradicts the claim that the same branch blocks smooth continuation at \(T_*\).
\end{proof}

\begin{theorem}[CM branch theorem]
Let \(O\) be a same-solution obstruction reached by the forward-positive Navier--Stokes program.
Then the pass/fail split is decided by the working class object:
\[
  O_{\Pass}\Longrightarrow \Member(Q;\Owork)\Longrightarrow\text{smooth continuation},
\]
and
\[
  O_{\Fail}\text{ capable of finite nonsmoothness}
  \Longrightarrow
  \Exit(Q;\Owork).
\]
\end{theorem}

\begin{proof}
The pass branch retains the carrier, participation, and field services on one witness record.
It is therefore a member branch and continues smoothly by membership readout.

The fail branch is offered as a finite same-solution terminal obstruction.
By finite-breakdown entry, it is CM-test-admissible.
By terminal packet capture, it loses Pack, or loses Part after Pack survives, or loses every positive Field scale after Pack and Part survive.
By finite failure witness class exit, that first failed service gives \(\Exit(Q;\Owork)\).
\end{proof}

\begin{theorem}[Clay closure from no genuine exit]
Assume the original smooth-data evolution has no genuine finite CM exit:
\[
  \text{OriginalSmoothData}(u_0)
  \Longrightarrow
  \neg\GenuineCMExit(T_*;u_0)
  \qquad\text{for every }T_*<\infty.
\]
Then the smooth Navier--Stokes solution generated by \(u_0\) has no finite terminal breakdown.
\end{theorem}

\begin{proof}
Assume, toward contradiction, that the maximal smooth solution from \(u_0\) stops at a finite time \(T_*\).
The alleged finite terminal object is CM-test-admissible by finite-breakdown entry.
The branch field has no third in-class terminal branch.

If the terminal branch is a member branch, membership readout extends the same solution past \(T_*\).
If the terminal branch is a fail branch, the CM branch theorem gives \(\Exit(Q;\Owork)\), so \(\GenuineCMExit(T_*;u_0)\) holds.
The first case contradicts maximality.
The second case contradicts the no-genuine-exit law.
Hence no finite terminal time exists.
\end{proof}

\section{The Short Proof Spine}

Once the definitions and branch law are in place, the proof spine is short.

\begin{claim}[Class proof spine]
For every obstruction \(O\) reached by the forward-positive Navier--Stokes program,
\[
  O_{\Pass}\Rightarrow\Member(Q;\Owork)\Rightarrow\text{smooth continuation},
\]
while
\[
  O_{\Fail}\Rightarrow\Exit(Q;\Owork).
\]
If original smooth data have no genuine finite CM exit, then no finite terminal breakdown remains.
\end{claim}

\begin{proof}
The first implication is membership readout.
The second implication is terminal packet capture followed by class exit.
The no-genuine-exit law removes the only fail branch that could be produced by the original smooth-data evolution.
The no-third-branch theorem removes an in-class nonsmooth terminal branch.
\end{proof}

\section{Referee Check}

The proof has three mathematical hinges.

First, terminal packet capture must be complete:
\[
  W_*
  \Rightarrow
  \neg\Pack_Q
  \vee
  \neg\Part_{N_s,Q}
  \vee
  \forall r>0\,\neg\Field_{N_s,r,Q}.
\]

Second, membership readout must really give the classical continuation bound:
\[
  \Member(Q;\Owork)
  \Longrightarrow
  \sup_{t<T_*}\|u(t)\|_{H^s}<\infty,\qquad s>5/2.
\]

Third, the original smooth-data evolution must satisfy the class-law closure:
\[
  \text{OriginalSmoothData}(u_0)
  \Rightarrow
  \neg\GenuineCMExit(T_*;u_0).
\]

A demand to prove one particular source, reserve, no-pulse, or readout estimate belongs to the forward-positive program.
Such a theorem can strengthen the pass side and simplify the branch field.
The class-membership contrapositive needs those estimates only when they attack one of the three hinges above.

\end{document}
